\subsection{Variational Quantum Eigensolver (VQE)}\label{sec:variational-quantum-eigensolver-vqe}

\subsubsection{The Physical Motivation}\label{sec:the-physical-motivation}

\textcolor{Red2}{\faIcon{exclamation-triangle} \textbf{The Fundamental Problem.}} Before introducing what the Variational Quantum Eigensolver (VQE) is, it is important to understand the physical motivation behind it. The \hl{VQE is used to find the \textbf{minimum-energy state} of a molecule by solving the Schrödinger equation.} Before understanding the mathematics, we first need to understand \textbf{why this problem is important}. Suppose we have a molecule. For example $\mathrm{H}_{2}$ or $\mathrm{H}_{2}\mathrm{O}$ or $\mathrm{CH}_{4}$. The first a chemist asks is: \emph{\textbf{is this molecule stable?}} And if it is stable, \emph{\textbf{what is its energy?}}

\highspace
\textcolor{Green3}{\faIcon{question-circle} \textbf{Why does energy matter?}} Nature has a very simple rule: systems naturally evolve toward lower energy states. For example, if you have a ball on a hill, it will roll down to the bottom of the hill. The same principle applies to molecules: electrons and nuclei arrange themselves until the molecule reaches its \textbf{lowest-energy configuration}. That configuration is called the \textbf{ground state}.

\highspace
\textcolor{Red2}{\faIcon{exclamation-triangle} \textbf{The Problem: Classical Difficulty.}} Almost every important property of a molecule depends on its ground state. For example, the color of a molecule, its reactivity, and its stability all depend on the ground state. Therefore, if we want to understand a molecule, we need to find its ground state.

\highspace
However, finding the ground state is \textbf{very difficult} for classical computers. The reason is that the number of possible configurations of electrons grows exponentially with the number of electrons in the molecule. For example, for a molecule with 10 electrons, there are $2^{10} = 1024$ possible configurations. For a molecule with 20 electrons, there are $2^{20} = 1,048,576$ possible configurations. For a molecule with 100 electrons, there are $2^{100} \approx 1.27 \times 10^{30}$ possible configurations! This exponential growth makes it impossible for classical computers to find the ground state of large molecules.

\highspace
\textcolor{Green3}{\faIcon{check-circle} \textbf{The Solution: Quantum Computers.}} A molecule is made of electrons and nuclei, which are \textbf{quantum particles}. Therefore, we cannot describe them using classical mechanics. We must solve quantum mechanics. Also, due to the exponential growth of the number of configurations, we need a quantum computer to solve this problem efficiently. Thus, this is exactly the motivation originally proposed by Richard Feynman: ``\emph{a quantum system should be simulated by another quantum system}''. Instead of storing exponentially many amplitudes on a classical computer, a \hl{quantum computer naturally evolves according to quantum mechanics.}

\highspace
\textcolor{Green3}{\faIcon{question-circle} \textbf{Where does VQE fit?}} Unfortunately, today's quantum computers are NISQ devices. So we cannot simply run a huge quantum simulation. Instead, the quantum computer is used to prepare candidate quantum states and classical computer is used to optimize the parameters of the quantum state. \textbf{Together they search for the lowest-energy state}. That \hl{hybrid algorithm is the Variational Quantum Eigensolver (VQE).}